\color{uniblue}\subsection[Выбор уставки токовой отсечки]{Выбор уставки токовой отсечки}
\color{black}

\begin{enumerate}

\item
По выражению (\ref{eq:to1}) определяется первичный ток срабатывания токовой отсечки $I_{\textsl{с.з.ТО}}$:
\begin{equation*}
    I_{\textsl{с.з.ТО}} = k_{\textsl{н}} \cdot I_{\textsl{к.макс.}}^{(3)} = 1,2 \cdot 581 = 697,2\; (\textsl{А})  ,
\end{equation*}

\item
Для оценки чувствительности токовой отсечки необходимо определить коэффициент чувствительности по формуле:
\begin{equation*}
k_{\textsl{ч}} = \frac{I_{\textsl{к.мин(k1)}}^{(2)}}{I_{\textsl{с.з.ТО}}}= \frac{0,87 \cdot 2570}{697,2} = 3,2 > 2,0,
\end{equation*}

\item
Уставка задается в относительных единицах от номинального тока аналогового входа устройства, для этого полученное значение пересчитывается по следующему выражению:
\begin{equation*}
    I_{\textsl{ср}} = \frac{I_{\textsl{с.з.МТЗ}}^I}{I_{\textsl{ном.вх.уст}}} \cdot \frac{I_{\textsl{ном.вт.ТТ}}}{I_{\textsl{ном.пер.ТТ}}} = \frac{697,2}{5} \cdot \frac{5}{600} = 1,16\; (\textsl{о.е.}),
\end{equation*}
\begin{eqexpl}[15mm]
    \item{$I_{\textsl{ном.вх.уст}}$} номинальный ток аналоговых входов устройства, 1 или 5 А;
    \item{$I_{\textsl{ном.вт.ТТ}}$} номинальный вторичный ток ТТ, А;
    \item{$I_{\textsl{ном.пер.ТТ}}$} номинальный первичный ток ТТ, А.
\end{eqexpl}
\item
Параметр \textbf{$I_{\textsl{ср}}$} <<Ток срабатывания>> функции ТО принимается \textbf{равным 1,16}.

\end{enumerate}